\documentclass[output=paper,colorlinks,citecolor=brown]{langscibook}
\ChapterDOI{10.5281/zenodo.20541862}

% \hyphenation{
%     We-li-nchen-kang-ci-kok
%     cross-ling-u-is-tic
% }

\author{Seongha Rhee\orcid{}\affiliation{Faculty of Liberal Arts, Mahidol University}}

\title{Grammaticalization of future markers in Korean}
\abstract{Contemporary Korean has a number of forms that make the future time reference. Their grammatical statuses are variable, e.g., an adnominalizing suffix, modal auxiliary, sentence-ender, etc. Unlike the cross-linguistically attested patterns, Korean futures involve two grammatical formants that inherently involve futurity, i.e., suffixes indicating ‘(at) a later time’ and ‘mode, manner, purpose’, which are primarily responsible for the emergence of future markers. In structural terms, the omissibility of sentential arguments and agglutinating morphology play an important role, because they make the grammatical status of a construction undergoing grammaticalization ambiguous and consequently promotes functional reinterpretation and structural reanalysis. A review of grammaticalization scenarios with respect to the grammaticalization parameters shows that grammaticalized future markers exhibit variable degrees of conformity to them.}

%move the following commands to the "local..." files of the master project when integrating this chapter

%\newcommand{}[1]{{\fontspec{Baekmuk Batang}{#1}}}

% \usepackage{multirow}
% \usepackage{fontspec}
%\usepackage{polyglossia}
%\setdefaultlanguage{english}
%\setotherlanguage{korean}



\IfFileExists{../localcommands.tex}{
   \addbibresource{../localbibliography.bib}
   \usepackage{tabularx,multicol}
\usepackage{url}
\urlstyle{same}

\usepackage[english]{babel}
 
\usepackage{langsci-optional}
\usepackage{langsci-lgr}
\usepackage{langsci-gb4e}

\usepackage{langsci-branding}

%\usepackage{tabularx}

% elena
\usepackage{graphicx} %package to manage images
\usepackage{paralist}

% sarah & lena
\usepackage[justification=centering]{caption}
\usepackage{rotating}

% robin
\usepackage{lscape}
\usepackage{tikz}
\usetikzlibrary{tikzmark}

% dominika
\usepackage{multirow}
\usepackage{hanging}

% olaf & dylan
\usepackage{subcaption}

% svetlana
\usepackage{verbatim}
\usepackage{enumitem}

% wiemer
%\usepackage[markup=underlined]{changes}
%\usepackage{hyphenat}

% wir
\usepackage{longtable}

% for crossref in the introduction jiabin
\usepackage{xr}
   % \newcommand*{\orcid}{}
%

\makeatletter
\let\thetitle\@title
\let\theauthor\@author
\makeatother

\newcommand{\togglepaper}[1][0]{
   \bibliography{../localbibliography}
   \papernote{\scriptsize\normalfont
     \theauthor.
     \titleTemp.
     To appear in:
     E. Di Tor \& Herr Rausgeberin (ed.).
     Booktitle in localcommands.tex.
     Berlin: Language Science Press. [preliminary page numbering]
   }
   \pagenumbering{roman}
   \setcounter{chapter}{#1}
   \addtocounter{chapter}{-1}
}
%
\newbool{bookcompile}
\booltrue{bookcompile}
\newcommand{\bookorchapter}[2]{\ifbool{bookcompile}{#1}{#2}}
%
% %\newcommand{\orcid}[1]{}
\graphicspath{ {./figures/} }
%
\let\eachwordone=\itshape


% Seongha
% \newcommand{\koreex}[1]{{\fontspec{Baekmuk Batang}{#1}}}

% Kutida
\newfontfamily\thaifont{Noto Sans Thai} 
\newcommand{\textthai}[1]{{\thaifont #1}} % für Thai


\newfontfamily\chinesefont{Noto Sans CJK SC}
\newcommand{\textchinese}[1]{{\chinesefont #1}} % Für Chinesisch

% wiemer
\def\hyph{-\penalty0\hskip0pt\relax}
% \definechangesauthor[color=NavyBlue]{SH}



\newcolumntype{K}{>{\raggedright\arraybackslash}p{0.6\textwidth}}  % Right column with 0.7 width


% \newfontfamily\krn[Scale=MatchLowercase,BoldFont=SourceHanSerif-Bold.otf]{SourceHanSerif-Regular.otf}
% \AdditionalFontImprint{Source Han Serif KO}
% \XeTeXlinebreaklocale 'ko'
% \renewcommand{\krn}{}



% Rhee
\newfontfamily\koreanfont{Noto Sans CJK KR} 
\newcommand{\textkorean}[1]{{\koreanfont #1}} % für Koreanisch

\renewcommand{\REF}[2][]{(\ref{#2}#1)}
   %% hyphenation points for line breaks
%% Normally, automatic hyphenation in LaTeX is very good
%% If a word is mis-hyphenated, add it to this file
%%
%% add information to TeX file before \begin{document} with:
%% %% hyphenation points for line breaks
%% Normally, automatic hyphenation in LaTeX is very good
%% If a word is mis-hyphenated, add it to this file
%%
%% add information to TeX file before \begin{document} with:
%% %% hyphenation points for line breaks
%% Normally, automatic hyphenation in LaTeX is very good
%% If a word is mis-hyphenated, add it to this file
%%
%% add information to TeX file before \begin{document} with:
%% \include{localhyphenation}

\hyphenation{
    con-tri-bute
    par-a-digm
    pe-ri-phra-ses
    achieve-ment
    accomp-lish-ment
    We-li-nchen-kang-ci-kok
    cross-ling-u-is-tic
    rec-ord-s
    trun-ca-ted
    de-ri-va-tion-al
    stem-de-ri-va-tion-al
    pre-sent
    in-flu-en-ced
    sys-tem
    ques-tions
    Bor-kovs-kij
    Da-ny-len-ko
    e-qui-va-lents
    re-cog-nized
    In-ter-fe-renz-schicht
    Aus-ei-nan-der-se-tz-ung
    Früh-neu-hoch-deutsch-kor-pus
    Ak-tions-art
    Smir-no-va
    Whi-le
    Ar-chive
    bi-ginnan
    Rei-chen-bach
    ger-ma-nis-ti-sche
    Re-fe-renz-kor-pus
    le-xi-ka-lisch-en
    Alt-hoch-deutsche
    Li-te-ra-tur
    dia-chron
    Dia-lek-to-lo-gie
}

\hyphenation{
    con-tri-bute
    par-a-digm
    pe-ri-phra-ses
    achieve-ment
    accomp-lish-ment
    We-li-nchen-kang-ci-kok
    cross-ling-u-is-tic
    rec-ord-s
    trun-ca-ted
    de-ri-va-tion-al
    stem-de-ri-va-tion-al
    pre-sent
    in-flu-en-ced
    sys-tem
    ques-tions
    Bor-kovs-kij
    Da-ny-len-ko
    e-qui-va-lents
    re-cog-nized
    In-ter-fe-renz-schicht
    Aus-ei-nan-der-se-tz-ung
    Früh-neu-hoch-deutsch-kor-pus
    Ak-tions-art
    Smir-no-va
    Whi-le
    Ar-chive
    bi-ginnan
    Rei-chen-bach
    ger-ma-nis-ti-sche
    Re-fe-renz-kor-pus
    le-xi-ka-lisch-en
    Alt-hoch-deutsche
    Li-te-ra-tur
    dia-chron
    Dia-lek-to-lo-gie
}

\hyphenation{
    con-tri-bute
    par-a-digm
    pe-ri-phra-ses
    achieve-ment
    accomp-lish-ment
    We-li-nchen-kang-ci-kok
    cross-ling-u-is-tic
    rec-ord-s
    trun-ca-ted
    de-ri-va-tion-al
    stem-de-ri-va-tion-al
    pre-sent
    in-flu-en-ced
    sys-tem
    ques-tions
    Bor-kovs-kij
    Da-ny-len-ko
    e-qui-va-lents
    re-cog-nized
    In-ter-fe-renz-schicht
    Aus-ei-nan-der-se-tz-ung
    Früh-neu-hoch-deutsch-kor-pus
    Ak-tions-art
    Smir-no-va
    Whi-le
    Ar-chive
    bi-ginnan
    Rei-chen-bach
    ger-ma-nis-ti-sche
    Re-fe-renz-kor-pus
    le-xi-ka-lisch-en
    Alt-hoch-deutsche
    Li-te-ra-tur
    dia-chron
    Dia-lek-to-lo-gie
}
   \boolfalse{bookcompile}
   \togglepaper[10]%%chapternumber
}{}

\begin{document}
\maketitle

\section{Introduction}

Whether Korean has future tense markers or not remains controversial. Early studies propose a tripartite system of distinguishing the past, the present, and the future (\cite{Joo1910}; \cite{Park1972}; \cite{Choe1937}; \cite{Hong1947}; \cite{Yoo1956}; among others), but later studies, notably from \citet{Ko1965} and onward (see also \cite{Nam1972}; \cite{Sohn1975}; \cite{Sung1976}; \cite{Suh1976}, \citeyear{Suh1977}; \cite{Mun2005}, \citeyear{Mun2009}; among others), tend to subscribe to the binary tense distinction (the past and the non-past), placing the future in the category of modality, based on the fact that the alleged future markers signal such modal notions as conjecture, inference, intention, etc. However, there is still a body of literature that subscribes to the tripartite tense distinction, e.g., Choi (\citeyear{Choi1995}, \citeyear{Choi1996}), Hong (\citeyear{Hong1987}, \citeyear{Hong2008}), Min (\citeyear{Min1991}, \citeyear{Min1998}), among others. 

The categorial nomenclature notwithstanding, Korean does have a number of markers that signal the ‘futurity’ of an event or state of affairs, hence ‘future markers’ from a broad perspective. These markers exhibit different historical depths, degrees of grammaticalization, levels of currency, and areas of specialization. This chapter addresses the development of those markers from a grammaticalization perspective. 

Korean is spoken by 81.7 million people \citep{Eberhard2023} mostly in and around the Korean peninsula. It was formerly regarded as an Altaic language as the eminent Finnish Altaic linguist Gustaf Ramstedt proposed in 1939 \citep{Ramstedt1997}, but in the modern linguistic research tradition, the Altaic hypothesis is contested (see \cite{Kim1992}; \cite{Sohn1999}; \cite{Song2005}, among others, for discussion of the issue). \textit{Ethnologue} (26\textsuperscript{th} ed., \cite{Eberhard2023}) lists Korean as a Koreanic language along with Jejueo, a language variety spoken in Jeju Island. Recently, it is frequently discussed under the geographically motivated classificatory term ‘Transeurasian languages’ \citep{Johanson2010}, often along with Japanese, since the two languages share many structural characteristics (\cite{NarrogRhee2013}; \cite{Narrog2018}; \cite{Robbeets2021}; see \cite{Tian2022} and \cite{Robbeets2022} for controversy). 
% Bei mehreren Autoren sind von Rhee in den Literaturangaben unterschlagen wurden. Jetzt bereits in bib-Datei hinzugefügt

The historical depth of texts written in \textit{Hangeul }(or \textit{Hankul}), the Korean alphabet, goes back to the 15\textsuperscript{th} century CE (1443). The invention of this alphabetic-featural script is a landmark event that divides Early Middle Korean (EMiK; the 10\textsuperscript{th} c. through the mid-15\textsuperscript{th} c.) and Late Middle Korean (LMiK; the mid-15\textsuperscript{th} c. through the 16\textsuperscript{th} c.), and a large body of texts has been electronically archived through government-led projects. Prior to the invention of \textit{Hangeul}, several different writing systems were invented using Chinese characters for their meaning (semantogram) or sound (phonogram) to represent Korean words or grammatical morphemes. 

\largerpage[0.5]

%Haspelmath 1995 steht nicht in der Literaturangabe von Rhee; nach der Recherche sind mehrere Veröffentlichungen von Haspelmath im Jahr 1995, welche?
Typologically, Korean is an agglutinating, head-final, SOV language. In line with the SOV word order characteristics in morphology, it is suffixal and postpositional. Multiple suffixes and postpositions may occupy slots in a relatively fixed order. Argument NPs are often omitted and NPs may occur without their postpositional particles, especially in speech, and interpretive flexibility thus created often affects or effects grammaticalization. In terms of verbal morphology, finite verbs are marked for tense, aspect, mood, and modality, whereas non-finite verbs are often marked with ‘converb’-markers (\textit{converbalia}, \cite{Ramstedt1903}), the suffixes whose main function is adverbializing a non-finite verb \citep[][]{haspelmath1995converb}. Finite verbs occurring sentence-finally are further marked by a sentence-type marker, such as declarative, imperative, interrogative, hortative, exclamative, etc., all of which forming a complex paradigm further modulated by the degree of honorification and politeness, commonly known as \textit{hwakyey }‘speech level’ (\cite{Rhee2020}).

This chapter is organized as follows: Section \ref{sec:2} addresses the development in form and function of the future markers; Section \ref{sec:3} discusses morphosyntactic and conceptual motivations and reviews the grammaticalization scenarios in view of grammaticalization parameters; and Section \ref{sec:4} summarizes the findings and concludes the chapter.


\section{Future time reference in Korean} \label{sec:2}

\largerpage[1]

In LMiK, the notion of future was carried by several different morphosyntactic devices, such as monosyllabic \textit{-li-} and the periphrastic forms involving the mode converb-marker -\textit{key} (e.g., \textit{-key taoy-, -key ha-,} and -\textit{kuys}-). However, these markers did not develop their tense-aspectual scopes, and they either survived in limited usages or completely disappeared, in Present-day Korean (PDK; the 21\textsuperscript{st} c.).

PDK has several future (and future-related) markers. Historically, they belong to three developmental trajectories, which form layers of different historical depths, i.e., \textit{-li-, -l.ke(s).i-,} and \textit{-keyss-.} Among these, the oldest form is \textit{-li-}, attested in Old Korean (OK; 1\textsuperscript{st} c. through early-10\textsuperscript{th} c.). On the other hand, the other two are more recent innovations; \textit{-l.ke(s).i-} developed from the 16\textsuperscript{th} century (\cite{Lee1997}: 64) and \textit{-keyss-} developed in the early 19\textsuperscript{th} century (\cite{Huh1987}; \cite{Rhee1996}). We discuss those three classes of future markers focusing on their historical development in form and function, in the following. 

\subsection{The \textit{-li-}Class futures}
\subsubsection{The future adnominalizer \textit{-l}} \label{sec:2.1.1}

The first future marker is \textit{-l}, a future adnominalizer, known as the ‘prospective adnominalizer’ (\textsc{pros.adn}) (\cite{Rhee2012b}), of which the major function is to transform a verbal or clausal constituent into a modifier of a noun phrase, hence the label ‘adnominalizer’. This function is illustrated in the following:


\ea \label{ex:rhee:1}
    \ea \label{ex:rhee:1a} \textkorean{올 사람}\\
    \gll \textit{o-l} \textit{salam}\\
     come-\textsc{fut} person\\
    \glt `a person who will come'

    \ex \label{ex:rhee:1b} \textkorean{살 옷}\\
    \gll \textit{sa-l} \textit{os}\\
     buy-\textsc{fut} clothes\\
    \glt `clothes that (I) will buy'
    \z
\z

%\il{Latin} %add "Latin" to language index for this page



Examples in \xref{ex:rhee:1} shows that \textit{-l} heads a constituent, which modifies the following noun (note that Korean is head-final), e.g. \textit{salam} ‘person’ in (1a) and \textit{os} ‘clothes’ in \xref{ex:rhee:1b}. Sometimes, structurally complex constituents involving the future \textit{-l} becomes fully unverbated, e.g., \textit{mek-ul kes }[eat-\textsc{fut} thing] ‘food’ (< ‘thing to eat’); \textit{ip-ul kes} [wear-\textsc{fut} thing] ‘clothing’ (note that \textit{-l} has a phonologically controlled allomorph -\textit{ul,} which occurs after a coda). As is obvious from the examples above, \textit{-l} is not a finite-verb morphology; it is used regardless of the speech time, as shown in \xref{ex:rhee:2}, in which \textit{-l} occurs in the sentences marked by the past tense. 


\ea \label{ex:rhee:2}
    \ea \textkorean{올 사람은 다 왔다}.\\
    \gll \textit{o-l} \textit{salam-un} \textit{ta} \textit{w-ass-ta}\\
         come-\textsc{fut} person-\textsc{top} all come-\textsc{pst-dec}\\
    \glt `Everyone who would come came.'\\
    \hfill (PDK, constructed)

    \ex \textkorean{그것은 내가 살 옷이었다}.\\
    \gll \textit{kukes-un} \textit{nay-ka} \textit{sa-l} \textit{os-i-ess-ta}\\
         it-\textsc{top} I-\textsc{nom} buy-\textsc{fut} clothes-\textsc{be-pst-dec}\\
    \glt `It was the clothes that I would buy.'\\
    \hfill (PDK, constructed)
    \z
\z

The examples in \xref{ex:rhee:2}, in which \textit{-l} occurs in the past-marked sentences, clearly shows that the future \textit{-l} is not a regular future-tense marker, which, by definition, is based on the speech time.\footnote{For this reason, \citet{Rhee2008} describes \textit{-l} ‘futuristic’ and exclusively uses the label ‘prospective adnominalizer’, avoiding ‘future’.} Thus, \textit{-l} is a prospective marker (as ‘prospective adnominalizer’, \cite{Rhee2008}), one belonging to the ‘relative tenses’ (\cite{McCawley1971}: 111; \cite{Partee1973}: 609; \cite{Comrie1985}; \cite{Dahl1985}; \cite{Michaelis2006}: 220; \cite{Hengeveld2011}), ‘secondary tenses’ (\cite{Diewald2011}; \cite{Lyons1977}: 689), ‘relational tenses’ (\cite{HopperTraugott2003}: 152), and ‘anaphoric tenses’ (\cite{LoCascio1986}: 198), in contrast with the absolute tenses anchored in the moment of speaking. Relative tenses are also called ‘compound tenses’ because these markers are, in general, morphosyntactically and/or semantically complex, by referring back to times denoted by previous temporal expressions (e.g., English \textit{will have done, had been doing}, etc.). In this regard, the Korean future \textit{-l} is qualitatively different in that it does not accompany a reference time expression. The future \textit{-l,} by virtue of marking a later point in time, may also be conceived as a ‘posterior’ tense (\cite{Jespersen1924}: 257; \cite{Bull1963}: 14; \cite{KampRohrer1983}: 254--257), as in the English posterior future ‘will be about to sing’ (\cite{Binnick2001}: 559).\footnote{The morpheme \textit{-l} is multifunctional, and an anonymous reviewer directs the author’s attention to \citet{Lim2009}, who lists the irrealis-marking usage among its diverse functions.}

Historically, the future-marking -\textit{l} had functioned as a nominalizer (e.g., gerundivizer, historically \textit{-l\textsc{h}, }with a glottal stop) with futuristic meaning (‘prospective nominalizer’), before it acquired the adnominalizer function (\cite{Kim1957}; \cite{lee1972koreanlinguistics}; \cite{Lee1997}; among others). This function survived as late as 1577 (\cite{Rhee2008}: 240--242) but has since become completely defunct; its sole function is a prospective adnominalizer (modifier) in Modern Korean. Its formal vestiges are found in certain sentence-enders, e.g., -\textit{lkka} (futuristic question), \textit{-lla} (apprehensional), -\textit{lseyla} (apprehensional), in which the futuristic morpheme \textit{-l-} is completely fused (cf. for discussion of interrogative markers involving \textit{-l} in OK and EMiK, see \cite{Ko2011}: 88--90, 101--127; for discussion on \textit{-l} in apprehensionals, see \cite{Yae2023}; \citeauthor{RheeKuteva2023}, in press). It has been noted that, crosslinguistically, the secondary grammaticalization from a nominalizer to a modifier (such as relativizer, complementizer, etc.) is fairly common (\cite{Yap2010}: 77; \cite{Yap2011}: 66). 

\largerpage[-0.5]

\subsubsection{The future \textit{-li-}}\label{sec:2.1.2}
The source of the future marker \textit{-li-} involves \textit{-l} (prospective nominalizer/ adnominalizer) addressed above. When \textit{-l} was combined with the defective noun \textit{i,} and the copula \textit{i-,} a new portmanteau form \textit{-li-} emerged (\cite{Lee1997}: 56--59; \cite{Park1998}: 163). As \textit{-l} had the prospective function, \textit{-li-} naturally inherited this function of marking futurity. Its use is exemplified in the following excerpts:\footnote{Some historical data in Korean has mixed script of Chinese characters and \textit{hangeul}, and a general practice in romanization of such texts is to use upper-case for Chinese characters following the modern Korean pronunciation. Since old Korean alphabets are typographically incompatible, Korean scripts are not presented if the excerpts contain such characters.}

\ea \label{ex:rhee:3}

    \ea \label{ex:rhee:3a}\textkorean{天心을 이루오리라}.\\
    \gll \textit{CHENSIM-ul} \textit{ilwu-o-li-la}\\
         providence-\textsc{acc} accomplish-\textsc{hon-fut-dec}\\
    \glt `I will fulfill the Providence (Heavens’ will).'\\
    \hfill (LMiK, 1445 \textit{Yongpiechenka} 108)


    \ex \label{ex:rhee:3b}\textkorean{내일이면 고향에 다다르리라}.\\
    \gll \textit{nayil-i-myen} \textit{kohyang-ey} \textit{tatalu-li-la}\\
         tomorrow-\textsc{be-cond} hometown-\textsc{at} reach-\textsc{fut-dec}\\
    \glt `I will arrive at (my) hometown by tomorrow.'\\
    \hfill (MoK, \textit{Wulimalsaym} ‘li’)


    \ex \label{ex:rhee:3c}\textkorean{나는 가리라, 곧 가리라 이니스프리섬으로}.\\
    \gll \textit{na-nun} \textit{ka-li-la,} \textit{kot} \textit{ka-li-la} \textit{inisuphuli.sem-ulo}\\
         I-\textsc{top} go-\textsc{fut-dec} soon go-\textsc{fut-dec} Innisfree.island-\textsc{dir}\\
    \glt `lit. I will go, To Innisfree Isle (I) will soon go.'\\
    \glt `org. I will arise and go now, and go to Innisfree.'\\
    \hfill (MoK translation, William B. Yeats, ‘The Lake Isle of Innisfree’)
    \z

\z

Examples in \xref{ex:rhee:3} show that \textit{-li-} is functioning as a future marker. In each example -\textit{li}- is followed by other verbal morphologies, e.g., a sentence ender, and, thus, the future marker \textit{-li-} is one that occupies a slot known as the penult ender position. Note that the clause containing the second occurrence of \textit{-li}- in \xref{ex:rhee:3c} is in the inverted word order for poetic effect. Even though \textit{-li-} was the primary future marker prior to MoK, it is now archaic or poetic. 

In terms of functions, the future \textit{-l-} signals intention or determination, as in \xref{ex:rhee:3a} and \xref{ex:rhee:3c} or conjecture as in \xref{ex:rhee:3b}. However, as functional classification largely depends on the variable granularity, other functions related to these overarching notions such as prediction, possibility, ability, obligation, etc. have been also marked by \textit{-li-.}\footnote{\citet[32--35]{Lee1997} also reports that in LMiK, the future \textit{-li-,} in combination with a negative adverb \textit{mot} ‘not’ or a negative verb \textit{mal}- ‘stop; do not’, marked prohibition, premonition, etc. (see \cite{Lee1997}: 35). }

\subsubsection{{The future sentence-ender \textit{-li}}}
The next future marker is the sentence-ender \textit{-li,} nearly identical in form and function with the previously addressed \textit{-li-} above. Its use is exemplified in the following: 

\ea \label{ex:rhee:4}

    \ea Declarative\\
    \gll \textit{YENGCWU-z} \textit{alph-\textsc{a}y} \textit{naynay} \textit{pwuskuli-li}\\
         wise.king-\textsc{gen} front-\textsc{at} forever be.ashamed-\textsc{fut.dec}\\
    \glt ‘(The unwise king) will be ashamed forever in the face of a wise king.’\\
    \hfill (LMiK, 1447 \textit{Yongpiechenka} 16)

    \ex Interrogative\\
    \gll \textit{CENGSACISEONG-skuy} \textit{nw-i} \textit{ani} \textit{o-z\textsc{a}v\textsc{a}-li}\\
         national.hero-\textsc{dat} who-\textsc{nom} not come-\textsc{hon-fut.q}\\
    \glt ‘Who would not follow this hero who defended the nation?’\\
    \hfill (LMiK, 1447 \textit{Yongpiechenka} 99)
    \z
\z


This future marker \textit{-li} is different from the previously addressed \textit{-li-} in that it is the sentence-ender (i.e., ultimate verbal morphology), whereas the latter is the one occurring in the penult verbal morphology. This is a unique situation because tense markers occur in the penult position and the ultimate position is reserved for sentence-type markers. The sentence-ending future markers occur only in declarative and rhetorical interrogative sentence types, as exemplified in \xref{ex:rhee:4} above. Since it is directly derived from the previous \textit{-li-,} the future sentence-ender \textit{-li} has the same functions, i.e., notably conjecture and intention, as well as other future-related functions.

This marker occurs very infrequently in the historical data, i.e., only in three sources that are written in the poetic style, i.e., \textit{Yongpiechenka} (1447), \textit{Welinchenkangcikok} (1447), and \textit{Welinsekpo} (1447) \citep[][20]{yeom1998sententialendings}. In MoK, the use of this future sentence-ender \textit{-li} is still restricted to poetic genres. For this reason, this future-marking sentence ender has not been discussed in earnest (\cite{Ko1990} and \cite{yeom1998sententialendings} are notable exceptions). The development of this unique marker results from a rather simple process. Sentence-final particles formerly attached to the tense-marking \textit{-li-} were omitted, presumably for poetic effects (e.g., to match the number of syllables with the counterpart stanza, cf. \cite{yeom1998sententialendings}:  20--21), and, consequently, the suffix was formally reanalyzed and functionally reinterpreted as a future-marking sentence-ender. For the lack of verbal morphologies that normally follow the tense marker, \citet[265]{Ko1990} and \citet[422]{Ryu1992} regard this ender as an instance of the \textit{panmal} ‘half-talk’ style, i.e., one without honorification or politeness markers. 


\subsection{The future \textit{-l.ke(s).i-}} \label{sec:2.2}
The second group of future markers consists of two historically related forms, \textit{-l.kes.i-} and \textit{-l.key}; the latter has diverged from the former. The future \textit{-l.kes.i-} developed from around the 16\textsuperscript{th} century (\cite{Lee1997}: 64) and is productively used in MoK.


\subsubsection{The future \textit{-l.kes.i-}}
The future tense marker \textit{-l.kes-i- }is a periphrasis involving the prospective (i.e., future-marking) adnominalizer \textit{-l} (see Section \ref{sec:2.1.1} above), the formal noun (i.e., semantically vacuous) \textit{kes} ‘thing’, and the copula \textit{i- }‘be’.\footnote{The lexical origin of the noun \textit{kes} is the noun \textit{kach/kas/kes} ‘skin, fur, surface, appearance, thing’ (\cite{Hong1983}; \cite{Rhee2008}: 248), which have been completely bleached. As an anonymous reviewer points out, the noun \textit{kes} in -\textit{l.kes.i-} may be replaced by other nouns, e.g., \textit{seym} ‘plan’, \textit{sayngkak} ‘thought’, etc., which form predicative constructions known as ‘mermaid’ constructions (\cite{Tsunoda2013}; \cite{Kim2013}; \cite{Rhee2020}: 598--599). Since these nouns still retain their lexical meaning, the mermaid constructions are regarded as syntactic constructions, unlike \textit{-l.kes.i-,} which is a relatively ‘pure’ future marker. The future meaning in the mermaid contructions is derived from the future adnominalizer -\textit{l,} discussed in Section \ref{sec:2.1.1}.} This complex construction has become fully univerbated through structural reanalysis and functional reinterpretation (see Section \ref{sec:3.1} below for more discussion). The future \textit{-l.kes.i-} is exemplified in the following: 

\ea
    \ea \textkorean{버스가 곧 올 것이다}.\\
    \gll \textit{pesu-ka} \textit{kot} \textit{o-l.kes.i-ta}\\
         bus-\textsc{nom} soon come-\textsc{fut-dec}\\
    \glt ‘The bus will come soon.’ \\
    \hfill (\cite{Lee2010}: 500)
    \z
\z

A brief note on the formal noun \textit{kes} is in order. \textit{Kes} is among the most frequently used nominalizers in MoK. It started its life as a full-fledged noun denoting ‘thing, skin, surface’ (\cite{Hong1983}; \cite{Rhee2011}) but became a defective noun from LMiK (\cite{Huh1983}, \cite{Jeong1989}; \cite{Lee1988}), labeled as such because it is so semantically weak that it cannot be used independently (also called ‘formal noun’). Even though it is semantically bleached, syntactically it still functions as the head noun of a nominalized string (\cite{Rhee2011}: 395), and thus requires that the modifying constituent be headed by an adnominalizer (such as the prospective adnominalizer \textit{-l} here). 

\largerpage[-1]

As the complex periphrasis \textit{-l.kes.i-} becomes univerbated into a single grammatical unit of signaling future, its function also develops into the conceptual domain marked by its earlier, and more strongly entrenched, competitor \textit{-li-.} Thus, \textit{-l.kes.i-} marks conjecture, intention, determination, possibility, prediction, obligation, etc. (see Section \ref{sec:2.1.2} above), largely depending on the context (in particular, the sentential subject) and the speech situation, as exemplified below (note that \textit{-ul} is an allomorph of \textit{-l}):

\ea \label{ex:rhee:7}
    \ea prediction, conjecture, conviction\\
    \textkorean{내일도 추울 것이다}.\\
    \gll \textit{nayil-to} \textit{chwuw-ul.kes.i-ta}\\
         tomorrow-\textsc{also} be.cold-\textsc{fut-dec}\\
    \glt ‘Tomorrow will be cold, too.’\\
    \hfill (\cite{Lee2010}: 500)\\

    \ex intention, determination\\
    \textkorean{나는 꼭 성공할 것이다}.\\
    \gll \textit{na-nun} \textit{kkok} \textit{sengkongha-l.kes.i-ta}\\
         I-\textsc{top} by.all.means succeed-\textsc{fut-dec}\\
    \glt ‘I will succeed by all means.’\\
    \hfill (\cite{Lee2010}: 501)\\
    \z
\z

In MoK, \textit{kes} [kʌs] itself has undergone phonological erosion and is often used in the form of \textit{ke} [kʌ] (note that \textit{-l.ke(s).i-} is only one of many other lexical and grammatical forms \textit{kes} participates in). Thus, the future marker \textit{-l.kes.i-} has its phonologically reduced variant \textit{-l.ke.i-, }from which the copula \textit{i- }may also be eroded due to the V-V sequence, which triggers deletion of one vowel or modification of the vowel quality. This process results in \textit{-l.ke-,} and thus the verb forms in \xref{ex:rhee:7} may also be \textit{chwuw-ul.ke-ta} `will be cold' and \textit{sengkongha-l.ke-ta} ‘will succeed’, respectively. There are no functional differences between these fuller and eroded variants, other than that the reduced variant is more closely associated with the colloquial styles.

\largerpage[1.5]

\subsubsection{The future \textit{-l.key}} \label{sec:2.2.2}

The future \textit{-l.key} originates from the reduced form of the future \textit{-l.kes.i-,} described above. Its usage is exemplified below: 

\ea \label{ex:rhee:8}
    \ea promissive\\ \label{ex:rhee:8a}
    \textkorean{내가 맛있는 것 사 줄게}.\\
    \gll \textit{nay-ka} \textit{masissnun} \textit{kes} \textit{sa-cwu-l.key}\\
         I-\textsc{nom} delicious thing buy-\textsc{ben-prom}\\
    \glt ‘I will buy something delicious for you.’\\
    \hfill (\cite{Lee2010}: 501)\\

    \ex polite command/request\\ \label{ex:rhee:8b}
    \textkorean{여기 침대에 누우실게요.}\\
    \gll \textit{yeki} \textit{chimtay-ey} \textit{nwuw-si-l.key-yo}\\
         here bed-\textsc{at} lie.down-\textsc{hon-imp-pol}\\
    \glt ‘Please lie on your back on the bed over here.’\\ \hfill (\cite{KooRhee2013}: 489)\\
    \z
\z

As shown in \xref{ex:rhee:8}, \textit{-l.key} is the sentence ender, i.e., occuring at the ultimate position (note that politeness marker \textit{-yo} is not considered in slot assignment), thus different from \textit{-l.kes.i-} occurring at the penult position, described in the previous section. 

When the reduced form is further marked by the ‘informal’ speech level, ‘universal’ speech-act sentence-ender \textit{-e/a,} the form contains a series of vowels as in -\textit{l.ke.i-a} /l.kʌ.i-a/, a situation highly susceptible to vowel change in Korean, thus resulting in -\textit{l.key} /l.ke/. This shows that the fusional process obliterates the word boundary as well as morpheme boundaries. The earliest attestations of \textit{-lkey} (e.g., 17\textsuperscript{th} c. \textit{Pakssipwuincen} movable type ed.) carried the same function as \textit{-l.kes.i-,} i.e., conjecture, prediction, intention, etc. (cf. \citep[490]{KooRhee2013}), but the reduced form gradually diverged, and from around the turn of the 20\textsuperscript{th} century, it became increasingly specializing as a promissive marker (‘I promise that...’) (e.g., 1911 \textit{Moktanhwa}; cf. Koo \& Rhee 2013: 490).

In the case of promissive, as in \xref{ex:rhee:8a}, since the speaker is also the agent of the proposed act, the promissive sentence is compatible with the first person subject only (either singular or plural). A Korean idiosyncrasy is that the understood subject may freely be omitted as is the case with other sentential constituents. Naturally, in promissive sentences, the implicit subject is the speaker, and the verb cannot be marked with honorification, because self-honorification is not acceptable in any situation. Therefore, when the predicate is marked with honorification, the understood sentential subject cannot be the speaker. When such forms are encountered in a sentence in which a putative promissive marker \textit{-l.key} is used, as in \xref{ex:rhee:8b}, in a situation where a nurse is giving an instruction to a patient, the addressee will make a pragmatic inference that since it cannot be the promissive sentence, it should be some other speech act. In other words, due to the omissibility of the sentential subject, the only signal that forces the addressee to interpret the sentence as an imperative is the presence of the honorification marker. Thus, the functional reinterpretation of the promissive as the imperative was actualized through pragmatic inference. Through this pragmatic inference, new functions of request (a milder form of imperative) or a command (a stronger form of imperative) emerged in PDK. 

\subsection{The future \textit{-keyss-}} \label{sec:2.3}
In PDK there are two primary futurity markers, \textit{-l.kes.i-} and -\textit{keyss}-. As shown in Section \ref{sec:2.2}, -\textit{l.kes.i-} etymologically involves the noun \textit{kes} ‘thing’. The latter, -\textit{keyss}-, which developed in the early 19\textsuperscript{th} century (\cite{Huh1987}; \cite{Rhee1996}), historically involves a number of morphemes (see below), especially, the mode converb-marker -\textit{key} ‘in such a manner that’, largely denoting mode, manner, purpose, etc. Its uses are illustrated in the following examples taken from \citet[67]{Lee2010} with modification:

\ea \label{ex:rhee:9}
\ea immediate future\\ \label{ex:rhee:9a}
\textkorean{잠시 후에 세 시가 되겠습니다}.\\
    \gll \textit{camsi} \textit{hwu-ey} \textit{sey} \textit{si-ka} \textit{toy-keyss-supnita}\\
         moment after-\textsc{at} 3 hour-\textsc{nom} become-\textsc{fut-dfr.dec}\\
    \glt ‘It will be 3 o’clock momentarily.’\\


\ex conjecture\\ \label{ex:rhee:9b}
\textkorean{어서 가자, 학교에 늦겠다}.\\
    \gll \textit{ese} \textit{ka-ca,} \textit{hakkyo-ey} \textit{nuc-keyss-ta}\\
         fast go-\textsc{hort} school-\textsc{to} be.late-\textsc{fut-dec}\\
    \glt ‘Let’s hurry up; we will/might be late for school.’\\


\ex possibility\\ \label{ex:rhee:9c}
\textkorean{나도 그만큼은 먹겠다}.\\
    \gll \textit{na-to} \textit{ku.mankhum-un} \textit{mek-keyss-ta}\\
         I-\textsc{also} that.much-\textsc{top} eat-\textsc{fut-dec}\\
    \glt ‘(I think) I, too, can eat that much. (It is possible that I...)’\\


\ex intention\\ \label{ex:rhee:9d}
\textkorean{내일은 미국으로 떠나겠어}.\\
    \gll \textit{nayil-un} \textit{mikwuk-ulo} \textit{ttena-keyss-e}\\
         tomorrow-\textsc{top} America-\textsc{dir} leave-\textsc{fut-end}\\
    \glt ‘I will leave for America tomorrow.’\\
    \z
\z

Example \xref{ex:rhee:9a} marks the immediate future, thus a type of the ‘metrical tense’ (\cite{ChungTimberlake1985}: 207ff; \cite{Comrie1985}: 83; \cite{Dahl1985}: 120ff; \cite{Binnick2001}: 560, among others). Even though the functions of \textit{-l.kes.i-} and -\textit{keyss}- overlap to a great extent, -\textit{keyss}- tends to specialize in immediate future (see Section \ref{sec:3.2} for discussion on its conceptual motivation). Example \xref{ex:rhee:9b} makes reference to conjecture as to the future event. Since conjecture can also make reference to the past event, -\textit{keyss}- is not restricted to future events (e.g., ‘They must have been late for school yesterday.’), in which case the future marker -\textit{keyss}- follows the past tense marker -\textit{ess}-. Example \xref{ex:rhee:9c} signals possibility but it may also mark its closely related concept ‘ability’, depending on the locus of the enabling force, i.e., within a sentient being (ability) or in the situation (possibility). Example \xref{ex:rhee:9d} marks the intention, typically the speaker’s in declarative and the addressee’s in interrogative. 

Even though the examples in \xref{ex:rhee:9} show the common functions of the future tense -\textit{keyss}-, there are a number of other functions that may be subsumed in those categories at varying degrees, e.g., ability, hypothetical willingness, current state as futuristic mode, evaluative, predestination, among others (\cite{Rhee1996}: 109--130).

The development of \textit{-keyss}- (and its dialectal form -\textit{kes(s)}-) has been studied by a number of researchers, from as early as \citet{Underwood1890}, \citet[]{Choe1989}, \citet[]{Park1972}, and \citet[]{Ramstedt1997} (see also \cite{Lee1947}; \cite{Na1953}; \cite{Shin1972}; \cite{Suh1977}; \cite{Gim1981}; \cite{Rhee1996}; \cite{Lee1997}, \citeyear{Lee2006}). The seemingly simple form -\textit{keyss-} originates from a construction involving multiple morphemes, i.e., the mode converb-marker -\textit{key} ‘in such a manner that’, the light verb \textit{ha}- ‘be\slash do\slash say\slash make’, the sequential converb-marker \textit{-e} ‘and’, the existence verb \textit{iss}- ‘exist’, even though many of them are not visible in the current form -\textit{keyss-}. The long process can be briefly described below.

The first is the emergence of the resultative/perfect/past marker -\textit{ess}-, developed from the converb-marker, also known as the ‘consolidating connective’ \textit{-e/a} (\cite{Koo1987}), and the verb of existence \textit{-is(s),} through a long process from the 15\textsuperscript{th} c. The consolidating connective, largely translatable as ‘and’, connects two sequentially related events and makes one holistic event (e.g., \textit{pel-e-mek-} ‘earn’ and ‘eat’ > ‘make a living’; \textit{kuw-e-mek-} ‘roast’ and ‘eat’ > ‘eat roasted meat’; etc.). By virtue of the presence of the existence verb in its historical origin, -\textit{ess}- tends to imply a strong relevance of a completed event to the present time (Note that it was a resultative marker before it developed into the perfective/past). 

The second is the emergence of the periphrastic causative \textit{-key.ha-,} developed from the mode converb-marker -\textit{key} and the light verb \textit{ha}- ‘do/be/say/make’. The converb-marker -\textit{key} is inherently futuristic because it refers to a resultant state of an action or event, e.g., \textit{aphu-key }[ache-\textsc{mod}] ‘in such a manner that one would feel pain; painfully’, \textit{cwuk-key} [die-\textsc{mod}] ‘in such a manner that one would die; leading to death’, `fatally' etc. The verb \textit{ha}- is a semantically bleached verb (thus the ‘light’ verb) with various general meanings. The combination of the ‘mode’ and ‘make’, thus ‘make x in such a manner that y will’, led to the development of the causative. 

The finally derived sense from prospective causality, definiteness, and current relevance is the standard notion of futurity, in that the speaker is making a judgment with respect to the future, based on the current definite situation which is a result of a completed event. Therefore, the notion `conjecture' is the most basic one in this marker. 


\section{Discussion} \label{sec:3}
In the preceding section, we have seen the development of various future markers, which can be categorized in three groups, one related to the prospective adnominalizer \textit{-l} (i.e., \textit{-l, -li-, }and -\textit{li}), one involving the defective nominal \textit{kes} ‘thing’ (i.e., \textit{-l.kes.i-} and \textit{-l.key}), and one involving the mode converb-marker -\textit{key} (i.e., -\textit{keyss}). Now we turn to a discussion of morphosyntactic motivation (Section \ref{sec:3.1}), conceptual motivation (Section \ref{sec:3.2}), and grammaticalization parameters (Section \ref{sec:3.3}).


\subsection{Morphosyntactic motivation} \label{sec:3.1}
There are a number of prominent factors of the development of the future in terms of morphosyntactic motivation. Among them are having a verb as the last component of the construction, the omissibility of arguments, and agglutinating morphology. 

We can see that the copula is involved in the development of most future markers, e.g., \textit{-li-, -li, -l.kes.i-, -l.ke.i-, }and \textit{-l.key.} In these markers, the last element \textit{i-} originates from the copula \textit{i-,} even though it is not directly visible in \textit{-l.key} (see Section \ref{sec:2.2.2}). The copula in the source constructions of these markers enables the preceding nominals/nominalizers (i.e., \textit{l} and \textit{kes}) to function as equative or identificational predicates by being a part thereof. For instance, we noted that the future marker -\textit{l.kes.i-} developed from a complex source construction involving the prospective adnominalizer \textit{-l}, the formal noun \textit{kes} ‘thing’, and the copula \textit{i- }‘be’, which became fully univerbated through structural reanalysis and functional reinterpretation. This is exemplified in the following:

\ea \textkorean{그가 올 것이다}. \label{ex:rhee:10}
    \ea source structure\\ \label{ex:rhee:10a}
    \gll \textit{ku-ka} \textit{o-l} \textit{kes-i-ta}\\
         he-\textsc{nom} come-\textsc{pros.adn} thing-\textsc{be-dec}\\
    \glt ‘lit. It is (a thing) that he comes at a later time.’\\

    \ex grammaticalized structure\\ \label{ex:rhee:10b}
    \gll \textit{ku-ka} \textit{o-l.kes.i-ta}\\
         he-\textsc{nom} come-\textsc{fut-dec}\\
    \glt ‘He will come.’\\
    \z
\z

The structural reanalysis of \xref{ex:rhee:10a} into \xref{ex:rhee:10b} involves a number of noteworthy aspects. As is obvious from the translation of \xref{ex:rhee:10a}, the entity being equated with ‘thing’ in the predicate is not the sentential subject ‘he’ (note that \textit{kes}‘thing’ is not used to refer to a person in Korean). The sentential subject is not explicit, i.e., the source structure of \xref{ex:rhee:10a} is a subject-less sentence, an idiosyncrasy of Korean syntax. The implicit subject is ‘the current situation’, i.e., the current situation is such that ‘he comes at a later time’. What aspects of the current situation warrants the conjecture are not explicit and thus unclear – they are only vaguely referred to as a ‘thing’, a perfect reference as its semantics is nearly vacuous (thus called a formal/defective noun, functioning as a nominalizer). Thus, the absence of concrete semantic content in \textit{kes} ‘thing’, the propinquity between the prospective adnominalizer \textit{-l} and \textit{kes}, the copula following \textit{kes}, and the presence of a contentful verb preceding all these components, among others, created a perfect syntagmatic condition for the reanalysis of the complex construction \textit{-l.kes.i-} as an auxiliary (note that the syntactically designated position of an auxiliary is between the main verb and the sentence-type marker like the declarative -\textit{ta} in \ref{ex:rhee:10}). 

It is also noteworthy that the reanalysis illustrated in \xref{ex:rhee:10} involves ‘syntactic upgrading’ (\cite{Heine1991}: 169--170; \cite{Ramm2008}; \cite{Rhee2012a}: 297--298), i.e., the formerly subordinating clause (‘he comes at a later time’) becomes the main clause with the future-tense interpretation (‘he will come’). Reanalysis of a former subordinate clause into a main clause typically occurs with ‘insubordination’ (\cite{Evans2007}), i.e., main clause ellipsis, which is responsible for many instances of grammaticalization in Korean, but the type of reanalysis with the future marker described here is actualized when the subject is implicit and its predicate is fused with the predicate of the subordinating clause. 

Another type of morphosyntactic motivation is with reference to the future -\textit{keyss}-. As alluded to in Section \ref{sec:2.3}, it developed from the string, -\textit{key-ha-e-iss-} [mode-make-and-exist], which developed into [\textsc{caus-perf}]. This is illustrated below:

\ea \textkorean{곧 비가 오게 하여 있어요}. \label{ex:rhee:11}
    \ea \label{ex:rhee:11a}
    \gll \textit{kot} \textit{pi-ka} \textit{o-key-ha-e-iss-e-yo}\\
         soon rain-\textsc{nom} come-\textsc{in.such.a.manner-make-and-exist-end-pol}\\
    \glt ‘It (the situation) is such a manner as to make the rain come soon and exists.’\\

    \ex \label{ex:rhee:11b}
    \gll \textit{kot} \textit{pi-ka} \textit{o-key.ha-ess-e-yo}\\
         soon rain-\textsc{nom} come-\textsc{caus-perf-end-pol}\\
    \glt ‘It has made the rain come soon.’\\

    \ex \label{ex:rhee:11c}
    \gll \textit{kot} \textit{pi-ka} \textit{o-keyss-e-yo}\\
         soon rain-\textsc{nom} come-\textsc{fut-end-pol}\\
    \glt ‘Rain will come soon.’\\
    \z
\z

One caveat is that the morphemic breakdown along with glosses and interpretations in the examples is grossly oversimplified, glossing over the complexity of the development; it is presented as such for the sake of convenience, hence the awkwardness of the interpretation. For instance, there is no instance in historical data that is comparable to \xref{ex:rhee:11a}, simply because the construction did not string together these exponents all at once at one point of time. In other words, there were always differential degrees of semantic cohesion and formal fusion between the components, which also changed over time in the course of grammaticalization (for detailed description of the processes, see \cite{Rhee1996}: 109--130 and \cite{Lee1997}: 64--82, among others). 

Simplicity of data presentation notwithstanding, the syntagmatic condition of -\textit{keyss}- is remarkably similar to that of \textit{-l.kes.i-,} in that the construction immediately follows a verb and is followed by a sentence-ender, a condition perfect for the exponent to be reanalyzed as an auxiliary. The implicit subject, one that takes the serialized verbs \textit{ha-} ‘make’ and \textit{iss}- ‘exist’ as its predicate, is something vague, such as ‘situation’, and does not surface in the sentence. As the two verbs became invisible due to formal fusion, the verb of the subordinating clause was reanalyzed as the main verb, and accordingly the fused exponent was reanalyzed as the auxiliary, once again an instance of syntactic upgrading. 

From the exposition above, it is evident that the syntactic idiosyncrasy in Korean that sentential constituents (the subject in the present case) can be freely omitted has played an important role in the grammaticalization of the future markers. In structural terms, the fact that the last element of the construction is a verb (e.g., \textit{i-} ‘be’ and \textit{iss-} ‘exist’ above) is important because the verb occurs sentence-finally, marked by the sentence-type marker (also known as the sentence-ender). The slot immediately before the sentence-type marker is for TAM, including the future auxiliary. Also significant is the typological characteristic of agglutinating morphology because agglutination and fusion obliterate the morphemic boundary, make the exponent’s grammatical status ambiguous, and consequently promotes functional reinterpretation and structural reanalysis.



\subsection{Conceptual motivation} \label{sec:3.2}
In crosslinguistic studies, the future is known to develop from diverse conceptual sources. For instance, the lexicon by \citet{Kuteva2019} lists Change-of-state, ‘come to’, ‘go to’, ‘love’, B-Necessity, D-Necessity, H-Possessive, ‘take’, ‘then’, ‘tomorrow’, Venitive, and ‘want’, as the commonly attested sources of future markers. In their seminal work, \citet[251--280]{BybeePerkinsPagliuca1994} show that lexemes related to desire (‘want’, ‘like’), ability (‘be able to’), obligation (‘owe’), movement verbs (‘come’, ‘go’), temporal adverbs (‘then’, ‘thereafter’, ‘afterward’, ‘just now’, ‘soon’), possession (‘have’, ‘get’, ‘obtain’, ‘catch’) as well as other semantically general verbs (e.g., ‘be’, ‘become’, ‘try to’, ‘need’, ‘owe’, ‘get’, ‘do’, ‘make’, ‘look for’, ‘intend’), etc. develop into the future across languages.

The Korean paths, however, are difficult to reconcile with these known crosslinguistic patterns. Even though the Korean verb \textit{ha}- ‘make, do’ may be close to the general verbs in the inventory, and the copula \textit{i-} ‘be’ is listed in it as well, we noted in the above that they do not make a notable contribution to the development. They are similar to syntactic placeholders rather than semantic contributors. Far more important contributors are the prospective adnominalizer \textit{-l} (e.g., ‘book to read’), which is already inherently futuristic by virtue of indexing a later point in time, and the mode converb-marker -\textit{key }(e.g., ‘painfully’), which inherently signals the effect- or goal-orientation (\cite{Rhee1996}). As briefly alluded to in Section \ref{sec:2.3}, the converb suffix -\textit{key} marks ‘manner’, ‘mode’ or ‘purpose’. The manner/mode of a situation or the purpose of an action has direct implication to the future because they are the bridge between one situation and its future situation or between one action and its future resultant situation. For this reason, the future signaled by -\textit{keyss}-, as compared to \textit{-l.kes.i-} or \textit{-li-,} is close to the present time, hence, the immediate future. Even though the relation between manner/purpose and future is almost universal, such relation implied by -\textit{key} in Korean is so prominent that \citep[][91]{Ramstedt1997} calls the verbs connected by the morpheme -\textit{key} ‘converbum futuri’, a position tantamount to classifying -\textit{key} as a full-fledged future marker. 

All the conceptual contributors to the development of Korean future markers can be shown as a diagram in Figure \ref{fig:rhee:1}, which is necessarily brief.

\begin{figure}
    \centering
    \includegraphics[width=0.8\linewidth]{figures/rhee/figure1.pdf}
    \caption{Conceptual contributors to Korean future markers}
    \label{fig:rhee:1}
\end{figure}

\largerpage[1]

\subsection{Grammaticalization parameters} \label{sec:3.3}
The grammaticalization scenarios of Korean future markers can be reviewed with respect to the four grammaticalization parameters suggested in \citet{Kuteva2019}, i.e., desemanticization, extension, decategorialization, and erosion. 


\subsubsection{Desemanticization}
Desemanticization refers to the loss of meaning content. Even though it is difficult to compare the source (or earlier) and target meanings, a rough generalization from the foregoing exposition can be presented as in Table \ref{tab:rhee:1}.


\begin{table}
\caption{Desemanticization in grammaticalization of future markers}
\label{tab:rhee:1}
\small
 \begin{tabularx}{.9\textwidth}{p{1.5cm}Xp{2.5cm}X}
  \lsptoprule
       {Form}     & {Source meaning\newline (literal)} & {Target meaning} & {Desemanticization}\\
  \midrule
{\textit{-l}} & ‘later time’ & Future & Not prominent \\\midrule
{\textit{-li-}} & ‘be (at) a later time’ & Future & Not prominent \\\midrule
{\textit{-li}} & ‘be (at) a later time’ & Future & Not prominent \\\midrule
{\textit{-l.ke(s).i-}} & ‘it is the thing that\newline P at a later time’ & Future & ‘thing’ disappeared \\\midrule
{\textit{-l.key}} & ‘it is the thing that\newline P at a later time’ & ‘I promise that...’ ‘Please do x’ & ‘thing’ disappeared \\\midrule
{\textit{-keyss-}} & ‘X is in the mode of\newline Y and exists’ & Future & `mode’ bleached,\newline `exist’ bleached \\
  \lspbottomrule
 \end{tabularx}
\end{table}



Table \ref{tab:rhee:1} shows that desemanticization in the \textit{-li-} Class futures is not prominent, because they started out (at least from the earliest available data) as a word/morpheme denoting a later time (either as a nominalizer or adnominalizer). The words denoting ‘thing’ (\textit{kes}) and ‘exist’ (\textit{iss}-) as well as the converb-marker with the ‘mode’ meaning (-\textit{key}) have their source meanings either completely absent or partially bleached. 


\subsubsection{Extension}
Extension as a grammaticalization parameter refers to the use of a form in new contexts. A comparison of the use contexts of the source structure to those of the current, future usage can be presented as Table \ref{tab:rhee:2}.




\begin{table}
\caption{Extension in grammaticalization of future markers}
\label{tab:rhee:2}
\scriptsize
 \begin{tabularx}{.97\textwidth}{p{1.3cm} p{3.8cm} p{3.7cm} p{1.6cm}}
  \lsptoprule
       {Form} & {Source context} & {Target context} & {Extension} \\
  \midrule
{\textit{-l}} & Verbal morphology\newline (Nominalizer) & Verbal morphology forming a\newline pre-nominal modifier & No extension (context shift) \\\midrule
{\textit{-li-}} & Verbal morphology for\newline predicate-formation (cf. \textit{-l-i-}) & Verbal morphology,\newline modal suffix & No extension (context shift) \\\midrule
{\textit{-li}} & Verbal morphology,\newline modal suffix & Verbal morphology,\newline sentence ender & No extension (context shift) \\\midrule
{\textit{-l.ke(s).i-}} & Verbal morphology involving\newline ‘a thing’ & Verbal morphology without\newline referential restriction & Extension \\\midrule
{\textit{-l.key}} & Verbal morphology involving\newline ‘a thing’ & Verbal morphology without\newline referential restriction & Extension \\\midrule
{\textit{-keyss-}} & Verbal morphology involving\newline causative and perfective meaning & Verbal morphology with\newline no causative\slash perfective meaning & Extension\slash Shift \\
  \lspbottomrule
 \end{tabularx}
\end{table}



As shown in Table \ref{tab:rhee:2}, the future markers involving the lexeme \textit{kes} ‘thing’ (i.e., \textit{-l.ke(s).i-} and \textit{-l.key}) have undergone extension in that the future does not have to involve reference to an entity. Similarly, -\textit{keyss}-, which contains elements denoting causative and perfective in earlier stages, can be used without such restrictive meanings. Extension, by definition, means an increase in use contexts and, thus, implies the availability of earlier usage at a later stage. However, the future markers involving \textit{-l} (i.e., \textit{-l, -li-,} and \textit{-li}) do not retain earlier usage, and thus exhibit usage context shift, rather than extension.

\subsubsection{Decategorialization}
Decategorialization refers to the loss in morphosyntactic properties characteristic of lexical or other less grammaticalized forms. Korean future markers display decategorialization at varying degrees, as shown in Table \ref{tab:rhee:3}.






\begin{table}
\caption{Decategorialization in grammaticalization of future markers}
\label{tab:rhee:3}
\footnotesize
     \begin{tabularx}{.9\textwidth}{p{1.3cm} p{2.2cm}p{2.2cm}X}
  \lsptoprule
       {Form} & {Source category\newline (feature)} & {Target category\newline (feature)} & {Decategorialization} \\
\midrule
{\textit{-l}} & Nominalizer,\newline verbal suffix & Adnominalizer,\newline verbal suffix & Decategorialization\newline Loss of argument function \\\midrule
{\textit{-li-}} & (periphrastic)\newline (cf. \textit{-l-i-}) & Modal auxiliary & Not determinable \\\midrule
{\textit{-li}} & Modal auxiliary & Modal suffix\newline (sentence ender) & Decategorialization\newline Loss of ability to be followed by a sentence-ender \\\midrule
{\textit{-l.ke(s).i-}} & (periphrastic) & Modal auxiliary & Not determinable \\\midrule
{\textit{-l.key}} & (periphrastic) & Modal suffix\newline (sentence ender) & Not determinable \\\midrule
{\textit{-keyss-}} & (periphrastic) & Modal auxiliary & Not determinable \\
  \lspbottomrule
 \end{tabularx}
\end{table}



The presence/absence or degrees of decategorialization should be based on a single, major (or non-minor) category lexeme. However, as shown in Table \ref{tab:rhee:3}, most future markers are periphrastic at the source stage, and, thus, decategorialization is largely indeterminable with these forms. Though trivial in extent, some level of decategorialization is observed with \textit{-l} and -\textit{li.} As we discussed in Section \ref{sec:2.1.1}, \textit{-l} was a nominalizer, which later became an adnominalizer. As a nominalizer it can form a nominal constituent, which can function as a verbal argument. As it becomes an adnominalizer (i.e., modifier), such a nominal function is lost, an instance of decategorialization. The sentence-ender \textit{-li}, occupying the ultimate slot, is a development from the modal auxiliary \textit{-li-,} occupying the penult slot. Since modal auxiliaries are suffixes that are further marked by sentence-enders (i.e., sentence-type markers), by becoming a sentence-ender, \textit{-li} has lost such earlier categorial function. 


\subsubsection{Erosion}

Erosion refers to the reduction of phonological volume of a grammaticalizing form. Korean future markers are eroded at varying degrees, as shown in Table \ref{tab:rhee:4}.



\begin{table}
\caption{Erosion in grammaticalization of future markers}
\label{tab:rhee:4}
\begin{tabularx}{.85\textwidth}{p{1.5cm} p{2cm}p{2.5cm}X}
  \lsptoprule
       {Form} & {Source form} & {Target form} & {Erosion} \\
\midrule
{\textit{-l}} & /l/ & /l/ & No erosion \\
\midrule
{\textit{-li-}} & /li/ & /li/ & No erosion \\
\midrule
{\textit{-li}} & /li/ & /li/ & No erosion \\
\midrule
{{\textit{-l.ke(s).i-}}} & {/l.kʌs.i/}  & /l.kʌs.i/ & No erosion\\
& & /l.kʌ.i/ & Eroded (moderate)\\
\midrule
{\textit{-l.key}} & /l.kʌs.i.a/ & /l.ke/ & Eroded (extensive) \\
\midrule
{{\textit{-keyss-}}} & {/ke.ha.ʌ.is/} & /kes/ & Eroded (extensive) \\
& & /kʌs/ (dialectal) & Eroded (extensive) \\
  \lspbottomrule
\end{tabularx}
\end{table}



It should be noted that there is a discrepancy between the character-based romanization and phonemic representation of the markers. From Table \ref{tab:rhee:4}, it is evident that the erosion is variable. The eroded forms are only those that involve the noun \textit{kes} ‘thing’ /kʌs/ and the mode-marking converb connective -\textit{key} /ke/. 

From the foregoing exposition, it is evident that the grammaticalization scenarios of Korean future markers exhibit differing degrees of conformity to the grammaticalization parameters, proposed by \citet{Kuteva2019}. Some of them are indeterminable because of multiple forms participating in the source construction.


\section{Summary and conclusion} \label{sec:4}

This chapter addressed the grammaticalization of Korean future markers. Contemporary Korean has six forms that make the future time reference, i.e., \textit{-l, -li-, -li, -l.ke(s).i-, -l.key}, and -\textit{keyss}-. Their grammatical statuses are variable from an adnominalizing suffix (\textit{-l}), to a modal auxiliary (-\textit{li-, -l.ke(s).i-, -keyss-),} and to sentence-ender (-\textit{li, -l.key}). 

Unlike the crosslinguistically attested patterns of grammaticalization sources, e.g., lexemes denoting desire, ability, obligation, movement, time, possession, among others, the development of Korean futures is largely due to two grammatical formants that inherently involve futurity, i.e., suffixes indicating ‘(at) a later time’ (-\textit{l}) and ‘mode, manner, purpose’ (-\textit{key}). 

In structural terms, a number of factors contribute to the development of future markers, but, most notably, omissibility of sentential arguments and agglutinating morphology play an important role, because they make the grammatical status of a construction undergoing grammaticalization ambiguous (i.e., through omission and fusion), and consequently promotes functional reinterpretation and structural reanalysis. 

A review of grammaticalization scenarios with respect to the grammaticalization parameters, i.e., desemanticization, extension, decategorialization, and erosion, shows that future markers exhibit variable degrees of conformity to them in the course of their grammaticalization.




\section*{Abbreviations}

\begin{tabularx}{.5\textwidth}{@{}lQ@{}}
\textsc{acc} & accusative \\
\textsc{adn} & adnominalizer \\
\textsc{ben} & benefactive \\
\textsc{cau} & causative \\
\textsc{cond} & conditional \\
\textsc{conv} & converb-marker \\
\textsc{cop} & copula \\
\textsc{dat} & dative \\
\textsc{dec} & declarative \\
\textsc{dfr} & deferential \\
\textsc{dir} & directional \\
\textsc{end} & sentence-ender \\
\textsc{fml} & formal \\
\textsc{fut} & future \\
\textsc{gen} & genitive \\
\textsc{hon} & honorific \\
\textsc{hort} & hortative \\
\end{tabularx}%
\begin{tabularx}{.5\textwidth}{@{}lQ@{}}
\textsc{imp} & imperative \\
\textsc{lmik} & Late Middle Korean \\
\textsc{mod} & Mode-converb \\
\textsc{mok} & Modern Korean \\
\textsc{n} & noun \\
\textsc{nom} & nominative \\
\textsc{nomz} & nominalizer \\
\textsc{pdk} & Present-Day Korean \\
\textsc{perf} & perfective \\
\textsc{pol} & polite \\
\textsc{prom} & promissive \\
\textsc{pros} & prospective \\
\textsc{pst} & past \\
\textsc{q} & question/interrogative \\
\textsc{tns} & tense \\
\textsc{top} & topic \\
\textsc{v} & verb \\
\end{tabularx}


\section*{Acknowledgements}

\largerpage[1]

Special thanks go to Prof. Dr. Stefan Hartmann and Dr. Lena Schnee for their kind support, to an anonymous reviewer for useful comments and suggestions, and to Prof. Dr. Shaun Manning for kindly reading and commenting on an earlier version of this paper. This research was supported by the research fund of the Faculty of Liberal Arts, Mahidol University, and by the Ministry of Education of the Republic of Korea and the National Research Foundation of Korea (NRF-2023S1A5A2A01081160).
%\section*{Contributions}
%John Doe contributed to conceptualization, methodology, and validation. 
%Jane Doe contributed to writing of the original draft, review, and editing.

\sloppy
\printbibliography[heading=subbibliography,notkeyword=this]
\end{document}

